% Explainable Expected Goal Models for Player Assessment -- DS340W
% Compile with pdfLaTeX on Overleaf (IEEEtran is built in).
\documentclass[conference]{IEEEtran}
\IEEEoverridecommandlockouts
\usepackage{cite}
\usepackage{amsmath,amssymb,amsfonts}
\usepackage{graphicx}
\usepackage{textcomp}
\usepackage{xcolor}

\begin{document}

\title{Explainable Expected Goal Models For Player Assessment: Extending and Improving Cavus Biecek}

% Add department, university, city/country and email on extra lines inside each block, e.g.
% \IEEEauthorblockA{\textit{Dept. name} \\ \textit{The Pennsylvania State University} \\ City, Country \\ email}
\author{
\IEEEauthorblockN{Will Bannon}
\and
\IEEEauthorblockN{Rishi Chada}
}

\maketitle

\begin{abstract}
Expected Goals (xG) is a metric used in the field of football (soccer) analytics that estimates the probability that a shot results in a goal. Because of the fact that goals are rare: clubs, broadcasters, analysts, and marketmakers use this xG metric to judge the quality of a chance and finishing more reliably than simply off goals. Our parent paper from Cavus and Biecek proposed explainable tree-based xG models and utilized aggregated profiles (AP) to assess shot probability at the player and team level. We use 316,000 shots from the top five major football leagues between 2014-15 and 2020-21 onward. The random forest of the original profiles scores an AUC of 0.985 on its training rows but 0.80 on shots it has not seen, predicts 2.4 expected goals for every goal scored, and places only 9\% of player seasons in its own 95\% confidence interval. Our cross-fitted models succeed in covering 96\% of player-seasons, and our logistic regression model achieves the best calibration of 24 benchmarked models. After 2021, every model, including Understat's proprietary model, over-predicts the highest quality shots by over 1.4 percentage points, indicating the context within the sport has drifted away from easy shots. Our results show that xG based player assessments require held-out evaluation, consistent calibrations, and retraining on occasion.
\end{abstract}

% Optional keywords:
% \begin{IEEEkeywords}
% expected goals, explainable artificial intelligence, model calibration
% \end{IEEEkeywords}

\section{Introduction}

Football is known to be one of the lowest scoring team games in sports. In the Understat shot data used in this study, only 33,715 of the 316,002 shots taken in the top 5 leagues between 2014-15 and 2020-21 seasons resulted in a goal; that's around 10.l n link\%. Due to goals being so rare, the result of a lol match or even the entirety of a players goal tally over an entire healthy season can carry high variance and a high degree of chance \cite{mead2023}. Expected goals (xG) addresses this by assigning every shot a probability of becoming a goal (from 0 to 1), estimated from the historical data of shots taken in similar situations \cite{anzer2021}. Summing the probabilities gives us the number of goals a player or team would be expected to score based upon the chances they created, which shows us the quality of those chances rather than how many of them actually converted into goals \cite{anzer2021}.

xG has moved on from its academic origins into everyday use in football context. It is shown in public broadcasts next to the score, clubs use it to recruit players and review their own players' performance, and it is used by market makers to price matches \cite{mead2023}. Most xG models estimate scoring probability based upon a small set of shot features like distance and angle to goal, body part used (shot type), and the shot situation \cite{mead2023}. There are models built on tracking data and positions of defenders and the goaltender, but these models are expensive and proprietary, so study of these models are limited in the academic field \cite{anzer2021}. The academic field primarily consists of models built on freely available event data. As these models have shifted from logistic regression to tree ensembles such as XGBoost and Random Forests, they become harder to interpret and explain. Explainable Artificial Intelligence (XAI) methods have been used to explain these models. Examples of these methods included partial dependence and ceteris-paribus profiles \cite{cavus2022} and Shapley additive explanations (SHAP) \cite{anzer2021}, aim to show how these models turn shot features into meaningful accurate prediction. This does matter when we use a model's output to judge individual players.

Cavus and Biecek \cite{cavus2022} brought these threads of XAI and AP together. Using 315,430 shots from the top five leagues in our selection, then trained CatBoost, XGBoost, LightGBM, and RF models with the forester AutoML package in R, on the original data and the versions class-balanced via under- and over-sampling with the ROSE method. Their over-sampled RF reported an area under the curve of ROC (AUC-ROC) of 0.985, which is again above the AUC of about 0.80 for the comparable models based on event data \cite{mead2023}. They proceeded to then aggregated ceteris-paribus profiles from this model to check how chance quality (shot quality) changes with distance and angle for individual players like Lionel Messi, Robert Lewandowski, and Burak Yilmaz in the 2020-21 season, as well as for teams like Schalke and Bayern Munich in a single Bundesliga match. Their study offered a practical template for explainable xG to assess players despite limitations.

To discuss limitations in brief, four features of that study limit how much its player assessments can be trusted. Firstly, its performance is not reflected in unseen shots as oversampling was applied before the data split to mitigate class imbalance and as a result half of the over sampled test rows were identical copies of training rows, which certainly inflates the reported results \cite{kapoor2023}. Re-evaluated on shots it had not seen, the RF's AUC falls to around 0.80. Secondly, xG is a probability that is summed across a large sample size, so it needs be calibrated: in essence, shots assigned a probability of 0.2 should score 20\% of the time. The original study does not assess calibration at all, even though resampling as a method of class balancing can be expected to distort predicted probability \cite{goorbergh2022}. Thirdly, the published profiles vary distance on fixed angle, and vice versa, even when both are computed from the same shot location. This results in between 30\% and 54\% of points the profiles assess being shots that cannot occur upon the dimensions of a professional football pitch, which forces the model to extrapolate based upon historical data. Their system of profiles does not measure uncertainty in a way where players can be easily benchmarked against one another. Additionally no out-of-time test is performed in the parent paper's model, which is helpful to measure how the sport itself evolves over time.

In our paper, we reproduce the models of Cavus and Biecek \cite{cavus2022} from their own published code and extend their study in three directions. We begin by comparing logistic regression, random forest, and XGBoost on calibration and discriminations, using the same fixed split of the paper's 2014-20 seasons into training, test, and untouched validation sets. We then tested how well our models trained on the era generalize to the 184,758 shots from 2021-22 onward to the present. We finally rebuilt the aggregated profiles so that each shot moves only only along paths that are physically possible, also adding match-level bootstrap bands and permutations tests for differences among the players and teams. For a fair comparison, we only use the seven shot features of the original study, so that our differences in results reflect fair benchmarking, evaluation, and modeling choices rather than being beaten by virtue of improved features and information. We find the original random forest overpredicts 2.3 expected goals for every goal scored and places only 9\% of their playe-seasons in their own 95\% interval, against the 96\% placements of our cross-fitted models. A logistic regression with spline terms for distance and angle, that are fitted separately via shot type, matches the best-calibrated model of the original study's 24 variants. After 2021, every model, including Understat's own proprietary model, over-predicts the highest quality chances of the final decile. Under our improved models, Lewandowski and Yilmaz outperformed their chances in 2020-21 and Messi performed in line with his, whereas the original model rated all three below par.

The remainder of this paper shall be organized as such. Section II reviews related works on xG models, explainable machine learning, the evaluation of probabilistic prediction, generalization out of time. Section III will describe data and evaluation design. Section IV will present the models, profiles, and metrics through various figures. Section V reports the results, Section VI shall be a discussion of the implications and limitations. Section VII shall conclude the paper.

\section{Related Work}

This section reviews three studies that represent a few core directions in xG research: models built on event data with added context \cite{mead2023}, models built on tracking data \cite{anzer2021}, and explainable models for player \cite{cavus2022}. We compare them across four criteria that matter for using xG to judge players: the data and features used, the models that are compared, how the models were evaluated, and how they were tested beyond the seasons used to train them

\subsection{Event-Data Expected Goals Model}

Early shot-quality models estimated scoring probability based upon where a shot was taken {[\,]}, and later event-data studies confirmed distance and angle to goal as the strongest conventionally collected predictors {[\,]}. Mead et al. \cite{mead2023} tested whether context beyond the actual shot improves such models. Using the single-season Wyscout event data from the 2017-2018 season of the top five leagues, joined with player values from Transfermarkt, Elo ratings from ClubElo, and attendance from FBref, they built 20 features across three groups: standard shot features, player and team ability, and psychological pressure (for example goal difference when shot was taken). They compared logistic regression, a multilayer perceptron, random forest, AdaBoost, and XGBoost all on a stratified 70/30 split, tuning them via a 10 fold cross validation and ranking them by log loss. Tunes XGBoost on all five leagues performed the best, with a test log loss of 0.2818, a Brier score of 0.079, and an AUC of 0.80, but logistic regression followed close behind at 0.2856 and was the best model for the La Liga and Ligue 1 leagues. Player value and goal difference were the most valuable added features. There was an additional analysis that found that xG predicted teams' future results better than shots or goals. However it is important to consider that this study only used a single season and didn't test calibration beyond a simple brier score. It's finding that logistic regression almost matched gradient boasting on event data also motivated the model comparison to be found in our study.

\subsection{Tracking-Data Expected Goals Models}

Event data records where a shot was taken but not the context surrounding it. Anzer and Bauer \cite{cavus2022} addressed this oversight with synchronized positional and event data for 105,627 shots from Bundesliga between 2013-15 and 2019-2020. Their features added shooter's speed, number of defenders in the line of the shot, goalkeepers position, and pressure on shooter. They compared gradient boosting, logistic regression, AdaBoost, and random forest on 60/20/20 training, validation, test split. Their final model combined separated gradient boosting submodels for open-play footed shots, headers, direct free kicks, with every penalty kick assigned the fixed value of 0.766. It reached an AUC of 0.824, against their base 0.772 built on event data alone, which shows how much the extra context adds. SHAP values ranked distance and angle as the most influence features, with goalkeepers distance from goal following along with pressure on shooter. In a workshop assistant coaches from Bundesliga agreed with the model's classification of results as deserved in more than 85\% of 102 matches. Unfortunately the data is proprietary so this model cannot be reproduced or applied to other leagues. Calibration was checked but only in aggregate, as 151.6 expected goals against 150 real goals in 1,357 shots from start of 2020-21, and the small sample is the study's only test in a later seasons.

\subsection{Explainable Expected Goals Models}

Cavus and Biecek \cite{mead2023} focused on making xG models interpretable for performance analysis. They used 315,430 Understat shots from the top five leagues between 2014-15 through 2020-21, describes by the same seven event features as our study: distance, and angle to goal, shot type, situation, the preceding action, the minute, and home or away. Using the forester AutoML package in R, they trained CatBoost, XGBoost, LightGBM, and random forest models on the original data and on data balanced by random under and over sampling from the ROSE package, and scored each on an 80/20 split. They reported AUC, Brier score, and log loss, with the over-samples random forest reaching an AUC of 0.985. Based upon this model they built aggregrated ceteris-paribus profiles, which average how the xG of every shot changes when one feature varies, to compare Messi? Lewandowski, and Yilmaz in 2020-21 and Schalke 04 and Bayern in a single match. but, the over -sampling was applied before forester split the data, the copies of the same shot leaked into the training and test sets, which is known to inflate results \cite{kapoor2023}, and balancing the classes by resampling distorts the predicted probabilities \cite{goorbergh2022}. The profiles vary distance and angle independently also, even when both are computed off the same shot location, so many of those profile points describe physically impossible shots. Due to this study's data and code being public, we can reproduce and correct it, so it's a natural baseline for our work.

\subsection{Summary and Research Gap}

To summarize, Mead et al. \cite{mead2023} advances our basis of xG through contextual features, Anzer and Bauer \cite{anzer2021} through tracking data, and Cavus and Biecek \cite{cavus2022} through XAI. Here are the gaps in what they show us: None of these studies assess calibration beyond summary measures, even though xG values are summed into player and team totals, and none test whether the models track across later seasons. Additionally the tracking data model can't be reproduced by us. We address these gaps in our project by reproducing Cavus and Biecek's model from their public code, evaluating calibration on held-out and cross-fitted predictions, testing upon seasons from 2021-22 into the present, and rebuilding the aggregated profiles with physically possible shot-pathing and uncertainty bands.

\begin{thebibliography}{00}
\bibitem{mead2023} J. Mead, A. O'Hare, and P. McMenemy, ``Expected goals in football: Improving model performance and demonstrating value,'' \textit{PLoS ONE}, vol. 18, no. 4, Art. no. e0282295, 2023, doi: 10.1371/journal.pone.0282295.
\bibitem{anzer2021} G. Anzer and P. Bauer, ``A goal scoring probability model for shots based on synchronized positional and event data in football (soccer),'' \textit{Front. Sports Act. Living}, vol. 3, Art. no. 624475, 2021, doi: 10.3389/fspor.2021.624475.
\bibitem{cavus2022} M. Cavus and P. Biecek, ``Explainable expected goal models for performance analysis in football analytics,'' in \textit{Proc. IEEE 9th Int. Conf. Data Sci. Adv. Anal. (DSAA)}, Shenzhen, China, 2022, pp. 1--9, doi: 10.1109/DSAA54385.2022.10032440.
\bibitem{kapoor2023} S. Kapoor and A. Narayanan, ``Leakage and the reproducibility crisis in machine-learning-based science,'' \textit{Patterns}, vol. 4, no. 9, Art. no. 100804, 2023, doi: 10.1016/j.patter.2023.100804.
\bibitem{goorbergh2022} R. van den Goorbergh, M. van Smeden, D. Timmerman, and B. Van Calster, ``The harm of class imbalance corrections for risk prediction models: Illustration and simulation using logistic regression,'' \textit{J. Amer. Med. Inform. Assoc.}, vol. 29, no. 9, pp. 1525--1534, 2022, doi: 10.1093/jamia/ocac093.
\end{thebibliography}

\end{document}
